% ============================================================
%  APUNTES UNIVERSITARIOS — DERECHO PROCESAL CIVIL Y COMERCIAL
%  Unidad 17: Modos Anormales de Terminación del Proceso
%  Autor: Isaac Edgar Camacho Ocampo
%  Ciclo Lectivo: 2024
% ============================================================

%% -------- CLASE DEL DOCUMENTO --------
\documentclass[12pt,oneside]{book}

%% -------- CODIFICACIÓN Y BABEL --------
\usepackage[utf8]{inputenc}
\usepackage[T1]{fontenc}
\usepackage[spanish,es-nodecimaldot]{babel}

%% -------- MÁRGENES --------
\usepackage[
    top=2.5cm,
    bottom=2.5cm,
    left=3cm,
    right=2.5cm,
    headheight=15pt
]{geometry}

%% -------- MATEMÁTICA --------
\usepackage{amsmath}
\usepackage{amssymb}
\usepackage{mathtools}

%% -------- IMÁGENES --------
\usepackage{graphicx}
\usepackage{float}
\usepackage{caption}
\usepackage{subcaption}

%% -------- TABLAS --------
\usepackage{booktabs}
\usepackage{array}
\usepackage{longtable}
\usepackage{tabularx}

%% -------- LISTAS --------
\usepackage{enumitem}

%% -------- COLORES Y CAJAS --------
\usepackage{xcolor}
\usepackage[most]{tcolorbox}

%% -------- ENCABEZADOS --------
\usepackage{fancyhdr}

%% -------- CONTROL TIPOGRÁFICO --------
\usepackage{ragged2e}

%% -------- HIPERVÍNCULOS --------

%% -------- RUTAS DE IMÁGENES --------
\graphicspath{
    {./img/}
    {./graficos/}
    {./figuras/}
}

%% -------- METADATOS (variables) --------
\newcommand{\universidad}{Universidad de Buenos Aires (UBA)}
\newcommand{\facultad}{Facultad de Derecho}
\newcommand{\carrera}{Derecho}
\newcommand{\materia}{Derecho Procesal Civil y Comercial}
\newcommand{\titulo}{Unidad 17 --- Modos Anormales de Terminación del Proceso}
\newcommand{\autor}{Isaac Edgar Camacho Ocampo}
\newcommand{\anio}{2024}

%% -------- CONFIGURACIÓN DE PÁRRAFOS --------
\setlength{\parindent}{0pt}
\setlength{\parskip}{0.5em}

%% -------- CONTROL DE VIUDAS Y HUÉRFANAS --------
\clubpenalty=10000
\widowpenalty=10000

%% -------- CONFIGURACIÓN DE LISTAS --------
\setlist[itemize]{topsep=0.3em, itemsep=0.2em}
\setlist[enumerate]{topsep=0.3em, itemsep=0.2em}

%% -------- ÍNDICE --------
\setcounter{tocdepth}{2}

%% -------- ENCABEZADOS Y PIES --------
\pagestyle{fancy}
\fancyhf{}
\fancyhead[L]{\nouppercase{\leftmark}}
\fancyhead[R]{\materia}
\fancyfoot[C]{\thepage}
\renewcommand{\headrulewidth}{0.4pt}
\renewcommand{\footrulewidth}{0pt}

\fancypagestyle{plain}{
    \fancyhf{}
    \fancyfoot[C]{\thepage}
    \renewcommand{\headrulewidth}{0pt}
}

%% -------- COMANDOS MATEMÁTICOS --------
\newcommand{\abs}[1]{\left|#1\right|}
\newcommand{\norm}[1]{\left\lVert#1\right\rVert}
\newcommand{\vect}[1]{\mathbf{#1}}
\newcommand{\deriv}[2]{\frac{d#1}{d#2}}
\newcommand{\pderiv}[2]{\frac{\partial#1}{\partial#2}}

%% -------- CAJAS ACADÉMICAS --------

% Definición
\newtcolorbox{definition}[1]{
    enhanced,
    breakable,
    title={Definición: #1},
    fonttitle=\bfseries,
    colback=blue!3,
    colframe=blue!50!black,
    boxrule=0.8pt,
    arc=2mm
}

% Importante
\newtcolorbox{important}{
    enhanced,
    breakable,
    title={Importante},
    fonttitle=\bfseries,
    colback=yellow!8,
    colframe=orange!70!black,
    boxrule=0.8pt,
    arc=2mm
}

% Ejemplo
\newtcolorbox{example}[1]{
    enhanced,
    breakable,
    title={Ejemplo: #1},
    fonttitle=\bfseries,
    colback=green!4,
    colframe=green!40!black,
    boxrule=0.8pt,
    arc=2mm
}

% Fórmula / Regla
\newtcolorbox{formula}[1]{
    enhanced,
    breakable,
    title={Fórmula / Regla: #1},
    fonttitle=\bfseries,
    colback=purple!4,
    colframe=purple!50!black,
    boxrule=0.8pt,
    arc=2mm
}

% Ejercicio
\newtcolorbox{exercise}[1]{
    enhanced,
    breakable,
    title={Ejercicio: #1},
    fonttitle=\bfseries,
    colback=gray!5,
    colframe=gray!60!black,
    boxrule=0.8pt,
    arc=2mm
}

% Nota
\newtcolorbox{note}{
    enhanced,
    breakable,
    title={Nota},
    fonttitle=\bfseries,
    colback=white,
    colframe=gray!50,
    boxrule=0.6pt,
    arc=2mm
}

%% -------- HIPERVÍNCULOS Y METADATOS PDF --------
\usepackage[
    colorlinks=true,
    linkcolor=blue!50!black,
    citecolor=green!40!black,
    urlcolor=blue!60!black,
    pdftitle={Apuntes de Derecho Procesal Civil y Comercial},
    pdfauthor={Isaac Edgar Camacho Ocampo},
    pdfsubject={Apuntes universitarios},
    bookmarks=true
]{hyperref}

% ============================================================
%  INICIO DEL DOCUMENTO
% ============================================================
\begin{document}

%% -------- PORTADA --------
\begin{titlepage}

\thispagestyle{empty}

\begin{center}

\vspace*{1cm}

{\large \textsc{\universidad}}

\vspace{0.2cm}

{\large \textsc{\facultad}}

\vspace{0.2cm}

{\large \carrera}

\vspace{3cm}

{\Huge \textbf{\materia}}

\vspace{1cm}

{\LARGE \titulo}

\vspace{0.8cm}

\textit{Comprender los institutos procesales es convertir la norma en herramienta de justicia.}

\vspace{1.2cm}

\rule{0.70\textwidth}{0.5pt}

\vspace{0.5cm}

{\large Apuntes de estudio}

\vspace{0.5cm}

\rule{0.70\textwidth}{0.5pt}

\vfill

{\large \autor}

\vspace{0.3cm}

{\large \anio}

\end{center}

\vfill

\begin{flushleft}
\textit{Este es un modesto aporte a los estudiantes de ``Derecho Procesal Civil y Comercial'', no pretende sustituir el material oficial de la materia.}
\end{flushleft}

\end{titlepage}

%% -------- ÍNDICE --------
\tableofcontents
\newpage

% ============================================================
%  CAPÍTULO 1: ENCUADRE GENERAL
% ============================================================
\chapter{Encuadre general: modos de terminación del proceso}

\section{El modo normal y los modos anormales}

El proceso judicial puede concluir de dos maneras: a través de su \textbf{modo normal} --- la sentencia definitiva, dictada tras el tránsito por las etapas introductoria, probatoria y de alegatos --- o a través de alguno de los \textbf{modos anormales} (también denominados \emph{modos extraordinarios} por parte de la doctrina) previstos en el Código Procesal Civil y Comercial de la Nación (CPCCN).

\begin{definition}{Modo anormal de terminación del proceso}
Todo aquel mecanismo que pone fin al proceso judicial por una vía distinta al dictado de la sentencia definitiva. El CPCCN regula cinco supuestos en los artículos 304 a 318: el desistimiento, el allanamiento, la transacción, la conciliación y la caducidad de la instancia.
\end{definition}

\begin{note}
La denominación ``modo anormal'' es la adoptada por el CPCCN. Parte de la doctrina prefiere hablar de ``modos extraordinarios'' de terminación del proceso, por oposición al modo ``ordinario'' o ``normal'' (la sentencia). La diferencia es meramente terminológica.
\end{note}

\section{Los cinco institutos del CPCCN}

Los cinco institutos previstos en los artículos 304 a 318 del CPCCN son:

\begin{enumerate}
    \item \textbf{El desistimiento} (arts. 304--306): renuncia o abandono de la instancia o del derecho por parte de quien ejerce la pretensión.
    \item \textbf{El allanamiento} (art. 307): sometimiento total o parcial del demandado a las pretensiones del actor.
    \item \textbf{La transacción} (art. 308): contrato bilateral con concesiones recíprocas que extingue obligaciones litigiosas o dudosas.
    \item \textbf{La conciliación} (art. 309): acuerdo intraprocesal celebrado ante el juez y homologado por éste.
    \item \textbf{La caducidad de la instancia} (arts. 310--318): extinción del proceso por inactividad de la parte que tiene la carga de impulsarlo, una vez vencidos los plazos legales.
\end{enumerate}

Cada instituto es analizado a continuación en sus requisitos, efectos procesales y diferencias con figuras afines.

% ============================================================
%  CAPÍTULO 2: EL DESISTIMIENTO
% ============================================================
\chapter{El desistimiento (arts. 304--306, CPCCN)}

\section{Concepto}

\begin{definition}{Desistimiento}
Acto jurídico procesal mediante el cual quien ejerce una pretensión renuncia o abandona esa posición procesal. Puede recaer sobre la instancia en curso (desistimiento del proceso) o sobre el derecho sustancial mismo (desistimiento del derecho o de la acción).
\end{definition}

En términos generales, desistir equivale a renunciar o abandonar. Su nota característica es que constituye un acto de quien \emph{ejerce} una pretensión --- no de quien se defiende.

\section{Legitimación: ¿quién puede desistir?}

Solo puede desistir quien articula una pretensión:

\begin{itemize}
    \item \textbf{El actor}, respecto de la demanda principal.
    \item \textbf{El reconviniente}, respecto de la reconvención planteada.
\end{itemize}

El demandado que únicamente ejerce defensas no tiene pretensión propia y, por consiguiente, no puede desistir.

\begin{important}
Para desistir válidamente se requiere \textbf{doble legitimación}:
\begin{enumerate}
    \item \textbf{Legitimación sustancial}: capacidad de derecho respecto de la pretensión en cuestión.
    \item \textbf{Legitimación procesal}: si quien desiste actúa mediante representante o apoderado, éste debe contar con poder \emph{especial y expreso} que lo faculte específicamente para desistir --- ya sea del proceso o del derecho. El poder general de representación no es suficiente.
\end{enumerate}
\end{important}

\section{Modalidades: desistimiento del proceso y del derecho}

El CPCCN distingue dos modalidades con efectos radicalmente distintos.

\subsection{Desistimiento del proceso}

El desistimiento del proceso supone la renuncia a continuar con la instancia en curso, pero \textbf{no afecta el derecho sustancial}. En consecuencia:

\begin{itemize}
    \item El actor puede iniciar un nuevo proceso en ejercicio del mismo derecho, siempre que la acción no se encuentre prescripta.
    \item No genera cosa juzgada material sobre el derecho.
\end{itemize}

Desde el punto de vista de su carácter unilateral o bilateral, la distinción es la siguiente:

\begin{itemize}
    \item \textbf{Antes de la notificación del traslado de la demanda}: el desistimiento del proceso es \emph{unilateral}; el actor puede desistir sin requerir la conformidad del demandado.
    \item \textbf{Después de la notificación del traslado de la demanda}: el desistimiento del proceso requiere la \emph{conformidad del demandado}. El fundamento radica en que, una vez notificado, el demandado ya está vinculado a la relación procesal y el desistimiento podría perjudicarlo si el actor quisiera reeditarlo en otro proceso.
\end{itemize}

\subsection{Desistimiento del derecho (o de la acción)}

El desistimiento del derecho importa la renuncia al derecho sustancial mismo --- no solo al proceso en curso. Sus características son:

\begin{itemize}
    \item Es \textbf{siempre unilateral}: no requiere la conformidad de la parte contraria en ningún caso.
    \item El fundamento de su carácter unilateral es claro: el actor renuncia a su propio derecho, lo cual beneficia al demandado --- quien queda liberado de ser demandado nuevamente por la misma pretensión --- y no lo perjudica.
    \item Produce efectos análogos a los de la \textbf{cosa juzgada material}: el derecho no puede hacerse valer en otro proceso distinto.
\end{itemize}

\subsection{Cuadro comparativo}

\begin{table}[H]
\centering
\begin{tabularx}{\textwidth}{
    >{\raggedright\arraybackslash}p{0.28\textwidth}
    >{\raggedright\arraybackslash}X
    >{\raggedright\arraybackslash}X
}
\toprule
\textbf{Aspecto} & \textbf{Desistimiento del proceso} & \textbf{Desistimiento del derecho / acción} \\
\midrule
¿Quién puede hacerlo? & Solo el actor o reconviniente & Solo el actor o reconviniente \\
\addlinespace
Conformidad de la contraparte & Necesaria si ya se notificó el traslado de demanda & No se requiere (siempre unilateral) \\
\addlinespace
Efecto sobre el derecho & No lo extingue; puede replantearse en otro proceso & Lo extingue definitivamente \\
\addlinespace
Cosa juzgada material & No genera cosa juzgada material sobre el derecho & Opera con efecto de cosa juzgada material \\
\addlinespace
Posibilidad de nuevo proceso & Sí, mientras el derecho no esté prescripto & No; el derecho fue renunciado \\
\bottomrule
\end{tabularx}
\caption{Diferencias entre el desistimiento del proceso y el desistimiento del derecho}
\end{table}

\section{Oportunidad para desistir}

El desistimiento solo puede operar \textbf{antes del dictado de la sentencia definitiva}. Una vez pronunciada la sentencia, el proceso ha concluido por su modo normal y el desistimiento carece de objeto: la pretensión ya fue resuelta mediante la sentencia y el derecho ya fue reconocido o rechazado.

% ============================================================
%  CAPÍTULO 3: EL ALLANAMIENTO
% ============================================================
\chapter{El allanamiento (art. 307, CPCCN)}

\section{Concepto}

\begin{definition}{Allanamiento}
Acto jurídico procesal mediante el cual el demandado --- o el reconvenido, respecto de la reconvención --- se somete total o parcialmente a las pretensiones de la parte actora.
\end{definition}

El allanamiento es el modo anormal que opera desde el polo pasivo de la relación procesal: le corresponde al \textbf{demandado} (a diferencia del desistimiento, que corresponde al actor o reconviniente).

\section{Naturaleza: el allanamiento no es un reconocimiento de hechos ni de derecho}

\begin{important}
El allanamiento \textbf{no importa} el reconocimiento de la verdad de los hechos alegados por el actor, ni el reconocimiento del derecho invocado. Es simplemente un acto voluntario del demandado que expresa su decisión de no oponerse al progreso de la pretensión. En consecuencia, los hechos reconocidos en el marco de un allanamiento \textbf{no pueden invocarse como verdaderos} en otro proceso, como si hubieran sido afirmados en una sentencia.
\end{important}

Esta distinción es esencial y se proyecta directamente sobre los efectos de cosa juzgada, que se analizan en la sección \ref{allanamiento_cj}.

\section{Características}

\begin{itemize}
    \item Es un acto \textbf{unilateral}: no requiere la conformidad del actor. El demandado puede allanarse sin que el actor deba aceptarlo.
    \item Puede ser \textbf{expreso} o \textbf{tácito}.
    \item Puede ser \textbf{total} o \textbf{parcial}.
\end{itemize}

\section{Modalidades: expreso y tácito}

\textbf{Allanamiento expreso}: la parte demandada manifiesta de modo contundente, sin lugar a dudas, su voluntad de someterse a las pretensiones del actor. La declaración debe ser inequívoca.

\textbf{Allanamiento tácito}: surge de una actitud del demandado que, aun sin manifestación expresa, resulta \emph{incontrovertible} en cuanto implica someterse a las pretensiones actoras. Dada la mayor incertidumbre que genera, se exige que la actitud sea absolutamente incompatible con cualquier otra interpretación.

\begin{example}{Allanamiento tácito claro}
El demandado, sin oponerse a la pretensión del actor, se presenta ante el tribunal y satisface íntegramente la pretensión mediante el pago. Aunque no manifestó expresamente ``me allano'', el efecto es incontrovertiblemente ese.
\end{example}

\begin{example}{Allanamiento tácito dudoso --- demanda de desalojo}
En una demanda de desalojo, la parte demandada se presenta y deposita en el tribunal las llaves del inmueble. La situación puede no ser incontrovertible: podrían existir otras cuestiones en discusión (p.ej., deuda de alquileres), de modo que la actitud no equivalga necesariamente a un allanamiento. Debe analizarse el caso concreto.
\end{example}

\section{Efectos según la combinación total/parcial y cumplimiento/incumplimiento}

\begin{table}[H]
\centering
\begin{tabularx}{\textwidth}{
    >{\raggedright\arraybackslash}p{0.30\textwidth}
    >{\raggedright\arraybackslash}X
    >{\raggedright\arraybackslash}X
}
\toprule
\textbf{Supuesto} & \textbf{Allanamiento + cumplimiento de la pretensión} & \textbf{Allanamiento sin cumplimiento} \\
\midrule
¿Qué hace el juez? & Dicta resolución interlocutoria y manda al archivo las actuaciones (art. 307 CPCCN) & Hace lugar a la pretensión y dicta sentencia definitiva para habilitar la ejecución \\
\addlinespace
Tipo de resolución & Resolución interlocutoria & Sentencia definitiva \\
\addlinespace
Etapa siguiente & Archivo del expediente & Ejecución de sentencia \\
\bottomrule
\end{tabularx}
\caption{Efectos del allanamiento según cumplimiento o no de la pretensión}
\end{table}

\section{La cosa juzgada en el allanamiento}
\label{allanamiento_cj}

La resolución que admite el allanamiento produce \textbf{cosa juzgada}, pero \emph{con un alcance limitado}:

\begin{itemize}
    \item La cosa juzgada opera respecto de la \textbf{pretensión}, no respecto de los hechos.
    \item Los hechos involucrados en el allanamiento no pueden invocarse como verdaderos en otro proceso, ya que no fueron objeto de reconocimiento ni de sentencia sobre su verdad.
\end{itemize}

Esto es coherente con la naturaleza del allanamiento: al no importar reconocimiento de hechos ni de derecho, la cosa juzgada resultante solo alcanza a la decisión sobre la pretensión puntual resuelta.

% ============================================================
%  CAPÍTULO 4: LA TRANSACCIÓN
% ============================================================
\chapter{La transacción (art. 308, CPCCN --- arts. 1641 y ss., CCyCN)}

\section{Marco normativo}

\begin{table}[H]
\centering
\begin{tabularx}{\textwidth}{
    >{\raggedright\arraybackslash}p{0.22\textwidth}
    >{\raggedright\arraybackslash}X
}
\toprule
\textbf{Norma} & \textbf{Texto} \\
\midrule
Art. 308, CPCCN & Las partes podrán hacer valer la transacción del derecho en litigio con la presentación del convenio o suscripción de acta ante el juez. Éste se limitará a examinar la concurrencia de los requisitos exigidos por la ley para la validez de la transacción y la homologará o no. En este último caso, continuarán los procedimientos del juicio. \\
\addlinespace
Art. 1641, CCyCN & La transacción es un contrato por el cual las partes, para evitar el litigio o ponerle fin, haciéndose concesiones recíprocas, extinguen obligaciones dudosas o litigiosas. \\
\bottomrule
\end{tabularx}
\caption{Marco normativo de la transacción}
\end{table}

\section{Concepto y naturaleza}

\begin{definition}{Transacción}
Contrato bilateral en virtud del cual las partes, haciéndose concesiones recíprocas, extinguen obligaciones litigiosas o dudosas (art. 1641, CCyCN). En el ámbito procesal, constituye un modo anormal de terminación del proceso que requiere presentación ante el juez y homologación judicial (art. 308, CPCCN).
\end{definition}

La transacción es un \textbf{acto bilateral}: requiere el acuerdo de \emph{ambas} partes. A diferencia del allanamiento o del desistimiento --- que son actos unilaterales --- ninguna de las partes puede imponer la transacción a la otra.

Puede ser \textbf{judicial} (celebrada en el marco de un proceso iniciado y presentada ante el juez) o \textbf{extrajudicial} (celebrada fuera del proceso, aunque igualmente puede recaer sobre derechos litigiosos).

\section{Requisitos formales}

\begin{formula}{Requisitos formales de la transacción judicial}
\begin{enumerate}
    \item Debe celebrarse \textbf{por escrito} (art. 1643, CCyCN).
    \item Debe \textbf{presentarse ante el juez} de la causa en la que tramitan los derechos litigiosos, acompañando el instrumento firmado por los interesados; o bien suscribirse directamente un \textbf{acta ante el juez}.
    \item El juez examinará la concurrencia de los requisitos legales y la \textbf{homologará o no}. Si no la homologa, el proceso continúa su curso.
\end{enumerate}
\end{formula}

\begin{note}
Conforme al art. 1643 del CCyCN, si la transacción recae sobre derechos litigiosos, solo es eficaz a partir de la presentación del instrumento firmado por los interesados ante el juez de la causa.
\end{note}

\section{Materias excluidas de la transacción}

El art. 1644 del CCyCN establece que no pueden ser objeto de transacción:

\begin{itemize}
    \item \textbf{Derechos en que esté comprometido el orden público}: por ejemplo, la nulidad de un matrimonio no puede ser materia de acuerdo entre las partes.
    \item \textbf{Derechos irrenunciables}: por ejemplo, los derechos laborales de los trabajadores.
    \item \textbf{Derechos sobre relaciones de familia o el estado de las personas}: por ejemplo, no es posible que las partes acuerden su divorcio mediante transacción ni que pacten el estado civil.
\end{itemize}

\textbf{Excepción}: sí pueden ser objeto de transacción los \emph{derechos patrimoniales derivados} de las relaciones de familia, tales como:

\begin{itemize}
    \item Las cuestiones sobre alimentos de hijos menores.
    \item La división de bienes conyugales bajo el régimen de comunidad.
\end{itemize}

% ============================================================
%  CAPÍTULO 5: LA CONCILIACIÓN
% ============================================================
\chapter{La conciliación (art. 309, CPCCN)}

\section{Marco normativo}

\begin{table}[H]
\centering
\begin{tabularx}{\textwidth}{
    >{\raggedright\arraybackslash}p{0.22\textwidth}
    >{\raggedright\arraybackslash}X
}
\toprule
\textbf{Norma} & \textbf{Texto} \\
\midrule
Art. 309, CPCCN & Los acuerdos conciliatorios celebrados por las partes ante el juez y homologados por éste tendrán autoridad de cosa juzgada. \\
\addlinespace
Art. 36, inc. 2°, CPCCN & En cualquier momento del proceso el juez puede convocar a las partes para alcanzar una conciliación. \\
\addlinespace
Art. 360, inc. 1°, CPCCN & En la audiencia preliminar, el juez intentará conciliar a las partes y acercarlas a los fines de alcanzar un acuerdo. \\
\bottomrule
\end{tabularx}
\caption{Marco normativo de la conciliación}
\end{table}

\section{Concepto}

\begin{definition}{Conciliación (intraprocesal)}
Acuerdo alcanzado por las partes dentro del proceso, ante la presencia y a instancia del juez, que pone fin al litigio y, una vez homologado por el magistrado, produce autoridad de cosa juzgada (art. 309, CPCCN).
\end{definition}

La conciliación intraprocesal ubica al \textbf{juez} como actor central: es el magistrado quien convoca a las partes --- en la audiencia preliminar (art. 360 CPCCN) o en cualquier momento del proceso (art. 36 CPCCN) --- para intentar acercarlas a un acuerdo.

\section{Efectos}

Cuando las partes alcanzan un acuerdo conciliatorio:

\begin{itemize}
    \item El proceso no termina por sentencia, sino por la realización del acuerdo.
    \item El acuerdo extingue las pretensiones antagónicas de las partes.
    \item Una vez \textbf{homologado por el juez}, produce el efecto de \textbf{cosa juzgada}.
\end{itemize}

\section{Distinción entre transacción y conciliación}

Ambos institutos comparten la nota de que conducen a un acuerdo que el juez puede homologar. Sin embargo, presentan diferencias esenciales:

\begin{table}[H]
\centering
\begin{tabularx}{\textwidth}{
    >{\raggedright\arraybackslash}p{0.25\textwidth}
    >{\raggedright\arraybackslash}X
    >{\raggedright\arraybackslash}X
}
\toprule
\textbf{Aspecto} & \textbf{Transacción} & \textbf{Conciliación} \\
\midrule
Iniciativa & Surge de la voluntad propia de las partes & A instancia o convocatoria del juez \\
\addlinespace
Ámbito & Judicial o extrajudicial; puede recaer sobre derechos litigiosos o dudosos & Siempre intraprocesal (dentro del proceso) \\
\addlinespace
Naturaleza jurídica & Contrato (arts. 1641 y ss., CCyCN) & Acuerdo procesal celebrado ante el magistrado \\
\addlinespace
Estado del derecho & Puede recaer sobre derechos dudosos o litigiosos & Solo sobre derechos ya litigiosos (proceso iniciado) \\
\addlinespace
Homologación & Puede o no ser homologado; el juez examina la validez & Se celebra ante el juez; produce cosa juzgada al homologarse \\
\bottomrule
\end{tabularx}
\caption{Diferencias entre transacción y conciliación}
\end{table}

% ============================================================
%  CAPÍTULO 6: LA CADUCIDAD DE LA INSTANCIA
% ============================================================
\chapter{La caducidad de la instancia (arts. 310--318, CPCCN)}

\section{Introducción y advertencia sobre derecho comparado provincial}

\begin{important}
El análisis que sigue se circunscribe al sistema del \textbf{CPCCN} (Código Procesal Civil y Comercial de la Nación), aplicable en el fuero federal y en la Ciudad Autónoma de Buenos Aires. La regulación de la caducidad de la instancia \textbf{varía entre provincias}: cada código procesal provincial puede prever plazos, supuestos de improcedencia y modos de operarse distintos. Al momento de aplicar el instituto en la práctica, es imprescindible remitirse al código procesal de la jurisdicción correspondiente.
\end{important}

La caducidad de la instancia es el instituto más extensamente regulado en el capítulo de modos anormales (arts. 310--318 CPCCN) y, a la vez, el que mayor diversidad de interpretaciones ha generado en la jurisprudencia.

\section{Concepto}

\begin{definition}{Caducidad de la instancia}
Modo anormal de terminación del proceso que opera cuando la parte que tiene la \textbf{carga de impulsar el procedimiento} incurre en inactividad procesal durante los plazos establecidos en el art. 310 del CPCCN. Configurados esos plazos, el tribunal puede declarar extinguida la instancia --- de oficio o a petición de la parte contraria.
\end{definition}

La nota esencial es la \textbf{inactividad procesal}: el proceso se paraliza por falta de impulso de quien tiene la carga de avanzarlo.

\section{Ámbito de aplicación: procesos dispositivos}

La caducidad de la instancia es aplicable a los \textbf{procesos dispositivos}, en los cuales la carga del impulso procesal recae sobre las partes. En este tipo de procesos, la inactividad de la parte que debe impulsar el avance --- una vez vencidos los plazos legales --- habilita la extinción de la instancia, ya sea de oficio o a petición de la contraria.

\section{Fundamento}

El instituto encuentra su justificación en dos razones concurrentes:

\begin{enumerate}
    \item \textbf{Interés público}: los procesos judiciales no deben quedar indefinidamente paralizados. La inactividad prolongada de quien tiene la carga de impulsar el proceso genera la presunción de \textbf{abandono} de la instancia.
    \item \textbf{Tutela del órgano jurisdiccional}: la relación procesal vincula no solo a las partes sino también al juez. El órgano jurisdiccional no puede quedar supeditado indefinidamente a los tiempos que las partes decidan manejar en sus procesos.
\end{enumerate}

\section{Distinción con la prescripción de la acción}

Ambos institutos se fundan en una presunción de abandono generada por la inactividad, pero son distintos en su naturaleza y efectos:

\begin{table}[H]
\centering
\begin{tabularx}{\textwidth}{
    >{\raggedright\arraybackslash}p{0.28\textwidth}
    >{\raggedright\arraybackslash}X
    >{\raggedright\arraybackslash}X
}
\toprule
\textbf{Aspecto} & \textbf{Prescripción de la acción} & \textbf{Caducidad de la instancia} \\
\midrule
Naturaleza & Instituto de derecho sustancial & Instituto de derecho procesal \\
\addlinespace
¿Qué se extingue? & El derecho / la acción & La instancia (el proceso en curso) \\
\addlinespace
¿Desde cuándo corre? & Desde que nació el derecho exigible & Desde el último acto impulsor del proceso \\
\addlinespace
¿Cancela la pretensión? & Sí (si prescribió, no puede reclamarse) & No (puede iniciarse nuevo proceso si el derecho no prescribió) \\
\addlinespace
Cómputo & Según cada norma sustancial & Días corridos (incluye inhábiles, salvo ferias judiciales) \\
\bottomrule
\end{tabularx}
\caption{Diferencias entre prescripción de la acción y caducidad de la instancia}
\end{table}

\begin{important}
La caducidad de la instancia \textbf{no cancela la pretensión}. El actor puede promover un nuevo proceso ejerciendo el mismo derecho, siempre que éste no se encuentre alcanzado por la prescripción. Confundir ambos institutos es uno de los errores más frecuentes en la práctica.
\end{important}

La caducidad de la instancia se abre con la \textbf{promoción de la demanda}, independientemente de que se haya notificado o no el traslado de la demanda al demandado.

\section{Plazos de caducidad --- Art. 310, CPCCN}

\begin{formula}{Plazos del art. 310, CPCCN}
Se producirá la caducidad de la instancia cuando no se instare su curso dentro de los siguientes plazos:

\begin{enumerate}
    \item \textbf{Seis (6) meses}, en primera instancia.
    \item \textbf{Tres (3) meses}, en segunda o ulterior instancia, y en cualquiera de las instancias en el juicio sumarísimo, en el juicio ejecutivo, en las ejecuciones especiales y en los incidentes.
    \item \textbf{El plazo en que opere la prescripción de la acción}, si fuese menor a los indicados precedentemente.
    \item \textbf{Un (1) mes}, en el incidente de caducidad de instancia.
\end{enumerate}
\end{formula}

\subsection{Aclaración sobre el inciso 4°: el incidente de caducidad de instancia}

El inciso 4° merece una explicación particular. Cuando la parte contraria acusa la caducidad de la instancia, se da traslado a la otra parte para que la conteste; luego el juez resuelve haciendo lugar o rechazando la caducidad. Este trámite constituye en sí mismo un \textbf{incidente}, y ese incidente de caducidad también tiene su propio plazo de caducidad: \textbf{un mes}.

\begin{example}{Funcionamiento del inciso 4° --- caducidad del incidente de caducidad}
La parte contraria acusa la caducidad de la instancia principal. El juez ordena el traslado a la otra parte. Si quien acusó la caducidad no notifica ese traslado (no envía la cédula electrónica correspondiente), y transcurre un mes sin que el incidente sea impulsado, la parte que debía ser notificada --- y que jamás lo fue --- puede a su vez acusar la caducidad del propio incidente de caducidad.
\end{example}

\section{Cómputo del plazo --- Art. 311, CPCCN}

\begin{formula}{Reglas de cómputo del art. 311, CPCCN}
Los plazos del art. 310 se computarán:
\begin{itemize}
    \item \textbf{Desde}: la fecha de la última petición de las partes, o resolución o actuación del juez, secretario u oficial primero, que tenga por efecto \textbf{impulsar el procedimiento}.
    \item \textbf{Durante}: los días inhábiles (incluyendo feriados), salvo los correspondientes a las ferias judiciales.
\end{itemize}
\end{formula}

\begin{important}
Los plazos de caducidad de instancia \textbf{corren en días corridos} (incluyendo días inhábiles y feriados), a diferencia de los plazos procesales ordinarios que se computan en días hábiles. La única excepción: no se computan los días correspondientes a las \textbf{ferias judiciales} (feria de verano --- enero --- y feria de invierno --- quince días en julio).
\end{important}

\subsection{El concepto de acto impulsor}

No cualquier presentación en el expediente constituye un acto impulsor del procedimiento. Para interrumpir el plazo de caducidad, el acto debe:

\begin{itemize}
    \item Tener por finalidad concreta \textbf{hacer avanzar el proceso} hacia las etapas subsiguientes.
    \item Significar un impulso procesal real, no meramente formal o accesorio.
\end{itemize}

\begin{example}{Acto que NO constituye impulso procesal}
Presentarse en el expediente únicamente para autorizar a un colaborador del estudio jurídico a visualizar el expediente no constituye un acto impulsor válido a los efectos del cómputo del plazo de caducidad.
\end{example}

El plazo se computa desde el día \emph{siguiente} al último acto impulsor válido.

\section{Litisconsorcio --- Art. 312, CPCCN}

En los supuestos de litisconsorcio, el impulso del procedimiento realizado por \textbf{uno de los litisconsortes beneficia a todos los demás}. El fundamento radica en la \textbf{indivisibilidad de la instancia}: no puede haber caducidad para un litisconsorte y no para otro. Si uno impulsa, ese impulso aprovecha a todos.

\section{Supuestos de improcedencia --- Art. 313, CPCCN}

El art. 313 del CPCCN establece taxativamente los casos en que la caducidad de la instancia \textbf{no procede}:

\begin{enumerate}
    \item \textbf{En los procedimientos de ejecución de sentencia}, salvo que se trate de incidentes que no guarden relación estricta con la ejecución procesal forzada propiamente dicha.
    \item \textbf{En los procesos sucesorios y, en general, en los voluntarios}, también salvo los incidentes que en ellos se susciten.
    \item \textbf{Cuando los procesos estuvieren pendientes de alguna resolución y la demora en dictarla fuere imputable al tribunal}: la inactividad debe poder imputarse a la parte, no al propio órgano jurisdiccional. Ejemplo: si ya se dictó el auto para sentencia y la sentencia no se pronunció por razones propias del tribunal, esa demora no puede perjudicar a la parte.
    \item \textbf{Si se hubiere llamado autos para sentencia}, salvo que se dispusiere prueba de oficio; en este último caso, cuando la producción de esa prueba dependiere de la actividad de las partes, la carga de impulsar el procedimiento existirá desde el momento en que éstas toman conocimiento de las medidas ordenadas.
\end{enumerate}

\section{Legitimación y modo de operarse --- Arts. 314, 315 y 316, CPCCN}

\subsection{Legitimación para pedir la declaración de caducidad (art. 315)}

La declaración de caducidad puede ser pedida:

\begin{itemize}
    \item En \textbf{primera instancia}: por el \textbf{demandado}.
    \item En \textbf{el incidente}: por la parte \emph{contraria} a quien lo hubiere promovido.
    \item En \textbf{el recurso}: por la \textbf{parte recurrida}.
\end{itemize}

La petición debe formularse \textbf{antes de consentir} cualquier actuación del tribunal o de la parte posterior al vencimiento del plazo legal. Si quien tiene derecho a acusar la caducidad consiente un acto impulsor posterior al vencimiento, pierde la posibilidad de invocar la caducidad por el período anterior.

La petición se sustanciará únicamente con un traslado a la parte contraria.

\subsection{Declaración de oficio (art. 316)}

La caducidad puede ser declarada \textbf{de oficio} por el tribunal, sin otro trámite que la comprobación de los plazos del art. 310, pero siempre \textbf{antes de que cualquiera de las partes impulse el procedimiento}.

\subsection{Efecto de la caducidad en segunda instancia sobre los recursos}

Cuando la caducidad opera en segunda instancia, tiene un efecto adicional importante: la declaración de caducidad importa el \textbf{desistimiento del recurso} interpuesto por el peticionario. Si ambas partes apelaron y se declara la caducidad de la segunda instancia respecto del recurso de una de ellas, la Cámara no tratará ninguno de los dos recursos: la sentencia de primera instancia quedará firme.

\section{Efectos de la caducidad --- Art. 318, CPCCN}

\begin{formula}{Efectos del art. 318, CPCCN}
\begin{itemize}
    \item \textbf{Caducidad en primera o única instancia}: no extingue la acción. El actor puede ejercerla en un nuevo juicio. Las pruebas producidas en el proceso caducado pueden hacerse valer en ese nuevo proceso.
    \item \textbf{Caducidad en instancias ulteriores}: acuerda fuerza de cosa juzgada a la resolución recurrida (la sentencia de primera instancia queda firme).
    \item \textbf{Relación entre instancia principal e incidentes}: la caducidad de la instancia principal comprende la reconvención y los incidentes. Pero la caducidad de un incidente \emph{no} afecta a la instancia principal.
\end{itemize}
\end{formula}

\subsection{Límite: la prescripción de la acción}

Si bien la caducidad en primera instancia no extingue la acción, el actor solo podrá iniciar un nuevo proceso si el \textbf{derecho no se encuentra prescripto}. Si entre la caducidad y la intención de reiniciar el proceso el plazo de prescripción de la acción ya venció, la pretensión no podrá ejercerse.

\subsection{Síntesis de efectos según la instancia}

\begin{table}[H]
\centering
\begin{tabularx}{\textwidth}{
    >{\raggedright\arraybackslash}p{0.30\textwidth}
    >{\raggedright\arraybackslash}X
}
\toprule
\textbf{Instancia en que opera} & \textbf{Efecto} \\
\midrule
Primera o única instancia & No extingue la acción; puede iniciarse nuevo proceso; las pruebas producidas pueden hacerse valer en él \\
\addlinespace
Instancias ulteriores (segunda, etc.) & Acuerda fuerza de cosa juzgada a la resolución recurrida (la sentencia de primera instancia queda firme) \\
\addlinespace
Instancia principal & Comprende la reconvención y los incidentes \\
\addlinespace
Incidente & No afecta a la instancia principal \\
\bottomrule
\end{tabularx}
\caption{Efectos de la caducidad según la instancia en que opera}
\end{table}

% ============================================================
%  CUADRO CORNELL
% ============================================================
\chapter*{Cuadro de repaso --- Método Cornell}
\addcontentsline{toc}{chapter}{Cuadro de repaso --- Método Cornell}
\markboth{Cuadro de repaso --- Método Cornell}{}

\vspace{0.5em}

\noindent\textbf{Unidad 17 --- Modos Anormales de Terminación del Proceso (CPCCN)}

\vspace{0.8em}

\begin{table}[H]
\centering
\begin{tabularx}{\textwidth}{
    >{\raggedright\arraybackslash}p{0.28\textwidth}
    >{\raggedright\arraybackslash}X
}
\toprule
\textbf{Preguntas / palabras clave} & \textbf{Notas / respuestas} \\
\midrule

\textbf{¿Qué es un modo anormal de terminación del proceso?} &
Todo aquel que no es la sentencia definitiva. El CPCCN regula cinco supuestos en los arts. 304--318: desistimiento, allanamiento, transacción, conciliación y caducidad de la instancia. Parte de la doctrina los denomina ``modos extraordinarios''. \\[0.6em]

\textbf{¿Quién puede desistir y de qué?} &
Solo quien articula una pretensión: el actor o el reconviniente. Puede desistir del \textbf{proceso} (no extingue el derecho; puede replantearse en otro proceso) o del \textbf{derecho / acción} (extingue el derecho con efecto de cosa juzgada material). El desistimiento del proceso requiere conformidad del demandado si ya fue notificado el traslado. \\[0.6em]

\textbf{¿Qué legitimación se requiere para desistir?} &
Doble legitimación: sustancial (capacidad sobre el derecho) y procesal (el representante debe tener poder especial y expreso; el poder general no alcanza). \\[0.6em]

\textbf{¿Qué es el allanamiento y a quién corresponde?} &
Acto jurídico procesal del demandado (o reconvenido) por el que se somete total o parcialmente a las pretensiones del actor. Es unilateral, puede ser expreso o tácito, y \textbf{no implica} reconocimiento de la verdad de los hechos ni del derecho invocado. \\[0.6em]

\textbf{¿Cómo opera el allanamiento con y sin cumplimiento de la pretensión?} &
\textbf{Con cumplimiento}: el juez dicta resolución interlocutoria y manda al archivo. \textbf{Sin cumplimiento}: dicta sentencia definitiva para habilitar la ejecución de sentencia. \\[0.6em]

\textbf{¿Qué cosa juzgada produce el allanamiento?} &
Cosa juzgada respecto de la \textbf{pretensión}, no respecto de los hechos (que no fueron reconocidos como verdaderos). Los hechos no pueden invocarse como reconocidos en otro proceso. \\[0.6em]

\textbf{¿Qué es la transacción y en qué difiere de la conciliación?} &
La transacción es un \textbf{contrato bilateral} en que las partes se hacen concesiones recíprocas para extinguir obligaciones litigiosas o dudosas (art. 1641 CCyCN; art. 308 CPCCN). Surge de la voluntad de las partes y puede ser judicial o extrajudicial. La conciliación, en cambio, surge de convocatoria del juez y es siempre intraprocesal (art. 309 CPCCN). \\[0.6em]

\textbf{¿Qué materias no pueden ser objeto de transacción?} &
Derechos que comprometan el orden público, derechos irrenunciables (ej. laborales), y derechos sobre relaciones de familia o estado de las personas. Excepción: derechos patrimoniales derivados de aquellas (alimentos, división de bienes conyugales). \\[0.6em]

\textbf{¿Qué es la caducidad de la instancia y cuál es su fundamento?} &
Extinción de la instancia por \textbf{inactividad procesal} de quien tiene la carga de impulsarla, transcurridos los plazos del art. 310 CPCCN. Fundamentos: interés público en no mantener procesos paralizados indefinidamente; tutela del órgano jurisdiccional; presunción de abandono de la instancia. \\[0.6em]

\textbf{¿Cuáles son los plazos del art. 310 CPCCN?} &
(1) \textbf{6 meses} en primera instancia. (2) \textbf{3 meses} en segunda/ulterior instancia, juicio sumarísimo, ejecutivo, ejecuciones especiales e incidentes. (3) El plazo de prescripción si es menor. (4) \textbf{1 mes} en el incidente de caducidad de instancia. \\[0.6em]

\textbf{¿Cómo se computan los plazos de caducidad?} &
Desde el último acto impulsor del procedimiento. Corren en \textbf{días corridos} (incluyendo días inhábiles y feriados). Solo se excluyen los días de ferias judiciales (enero y julio). Diferencia esencial respecto de los plazos procesales ordinarios, que se cuentan en días hábiles. \\[0.6em]

\textbf{¿Cuándo no procede la caducidad (art. 313)?} &
En ejecución de sentencia (salvo incidentes ajenos a la ejecución forzada); en procesos sucesorios y voluntarios (salvo incidentes); cuando la demora sea imputable al tribunal; cuando se haya llamado autos para sentencia (salvo prueba de oficio a cargo de las partes). \\[0.6em]

\textbf{¿La caducidad extingue la acción?} &
\textbf{No} en primera instancia: puede iniciarse nuevo proceso y hacerse valer las pruebas producidas, siempre que el derecho no esté prescripto. \textbf{Sí} en instancias ulteriores: acuerda fuerza de cosa juzgada a la resolución recurrida (la sentencia de primera instancia queda firme). \\[0.6em]

\textbf{¿Qué relación hay entre caducidad de la instancia principal e incidentes?} &
La caducidad de la instancia principal \textbf{comprende} la reconvención y los incidentes. La caducidad de un incidente \textbf{no afecta} a la instancia principal. \\[0.6em]

\midrule

\multicolumn{2}{p{\textwidth}}{
\textbf{Resumen:} La Unidad 17 aborda los cinco modos anormales de terminación del proceso previstos en los arts. 304 a 318 del CPCCN, por oposición al modo normal (sentencia definitiva). El \textbf{desistimiento} --- del proceso o del derecho --- es un acto del actor o reconviniente: el desistimiento del derecho extingue la pretensión con efecto de cosa juzgada, mientras que el del proceso solo extingue la instancia sin afectar el derecho. El \textbf{allanamiento} es el acto unilateral del demandado que se somete a las pretensiones del actor sin reconocer hechos ni derecho, con cosa juzgada limitada a la pretensión. La \textbf{transacción} (contrato bilateral con concesiones recíprocas, arts. 1641 CCyCN y 308 CPCCN) y la \textbf{conciliación} (acuerdo intraprocesal a instancia del juez, art. 309 CPCCN) difieren en la iniciativa y el ámbito. La \textbf{caducidad de la instancia} (arts. 310--318 CPCCN) opera por inactividad procesal en los plazos de 6 meses (primera instancia) o 3 meses (segunda instancia, sumarísimo, ejecutivo e incidentes), computados en días corridos con exclusión de ferias judiciales; no extingue la acción en primera instancia pero sí en instancias ulteriores, donde acuerda fuerza de cosa juzgada a la resolución recurrida.
} \\

\bottomrule
\end{tabularx}
\end{table}

\end{document}
